\underline{Exercise 12.2}
\begin{enumerate}
% Question 1
\item[1.] Use the Lagrange-multiplier method to find the stationary values of $z$.
\begin{enumerate}
\item[(a)] $z = xy$, subject to $x + 2y = 2$
\item[(b)] $z = x(y + 4)$, subject to $x + y = 8$
\item[(c)] $z = x - 3y - xy$, subject to $x + y = 6$
\item[(d)] $z = 7 - y + x^{2}$, subject to $x + y = 0$
\end{enumerate}

% Question 2
\item[2.] In Prob.~1, find whether a slight relaxation of the constraint will increase or decrease the optimal value of $z$. At what rate?

% Question 3
\item[3.] Write the Lagrangian function and the first-order condition for stationary values (without solving the equations) for each of the following:
\begin{enumerate}
\item[(a)] $z = x + 2y + 3w + xy - yw$, subject to $x + y + 2w = 10$
\item[(b)] $z = x^{2} + 2xy + yw^{2}$, subject to $2x + y + w^{2} = 24$ and $x + w = 8$
\end{enumerate}

% Question 4
\item[4.] If, instead of $g(x, y) = c$, the constraint is written in the form of $G(x, y) = 0$, how should the Lagrangian function and the first-order condition be modified as a consequence?

\end{enumerate}
